% BHU PYQ -- Manually transcribed & error-corrected
% Paper: MTB-604 : Mechanics
% Exam: B.A./B.Sc. (Hons) Semester VI Examination, 2013-14

\documentclass[11pt,a4paper]{article}
\usepackage[a4paper,top=2.5cm,bottom=2.5cm,left=2.8cm,right=2.8cm,headheight=14pt]{geometry}
\usepackage{amsmath,amssymb}
\usepackage[T1]{fontenc}\usepackage[utf8]{inputenc}
\usepackage{lmodern,microtype,fancyhdr,enumitem,hyperref,lastpage,parskip}

\hypersetup{colorlinks=true,urlcolor=black,linkcolor=black,
  pdftitle={MTB-604: Mechanics -- BHU B.A./B.Sc. Sem VI 2013-14}}
\pagestyle{fancy}\fancyhf{}
\renewcommand{\headrulewidth}{0.4pt}\renewcommand{\footrulewidth}{0.4pt}
\fancyhead[L]{\small\textit{Banaras Hindu University}}
\fancyhead[R]{\small\textit{MTB-604: Mechanics}}
\fancyfoot[C]{\small Page \thepage\ of \pageref{LastPage}}

\newlist{parts}{enumerate}{2}
\setlist[parts,1]{label=(\alph*),leftmargin=*,itemsep=5pt,topsep=3pt}
\setlist[parts,2]{label=(\roman*),leftmargin=*,itemsep=3pt,topsep=2pt}
\newcommand{\pts}[1]{\hfill\textit{[#1]}}

\begin{document}
\begin{center}
  {\large\bfseries BANARAS HINDU UNIVERSITY}\\[2pt]
  {\normalsize B.A./B.Sc. (Hons) --- Semester VI Examination, 2013-14}\\[6pt]
  \rule{0.8\linewidth}{0.5pt}\\[6pt]
  {\large\bfseries Paper No. MTB-604: Mechanics}\\[4pt]
  {\normalsize Time: 3 Hours \hfill Full Marks: 70}
\end{center}
\rule{\linewidth}{0.4pt}
\smallskip
\noindent\textit{Instructions: Answer \textbf{five} questions including Question~1 which is compulsory. Figures in right-hand margin indicate marks.}
\bigskip

\noindent\textbf{Question 1.}
\begin{parts}
  \item Define the central axis of a given system of forces acting on a rigid body and show that the moment of the resultant couple about the central axis is less than the moment of the resultant couple about any point not on the central axis. \pts{3}
  \item A uniform beam of thickness $2b$ rests symmetrically on a perfectly rough horizontal cylinder of radius $a$. Discuss the stability of the equilibrium. \pts{3}
  \item Find the moment of inertia of an area bounded by the sector of a circular plate making an angle $2\alpha$ at the centre, about a line in the plane of the sector, perpendicular to its axis of symmetry and passing through the centre. \pts{3}
  \item Define kinetically equivalent mechanical systems. \pts{2}
  \item Show that the centres of suspension and oscillation of a compound pendulum are convertible (interchangeable). \pts{3}
  \item Define the centre of percussion for the motion of a body about a fixed axis. Find the position of the centre of percussion of a uniform rod suspended freely from one end. \pts{3}
  \item Write the expression for the moment of momentum about a fixed origin of a body moving in two dimensions. Explain the meaning of the symbols used. \pts{3}
  \item State the principle of conservation of moment of momentum of a rigid body moving under the action of finite forces. \pts{3}
\end{parts}

\medskip\rule{\linewidth}{0.2pt}

\medskip\noindent\textbf{Question 2.}
\begin{parts}
  \item Describe the invariants for any given system of forces acting on a rigid body and find the condition for the system to compound into a single force. \pts{6}
  \item Two forces $P$ and $Q$ act along the lines $y = x\tan\alpha$, $z = c$ and $y = -x\tan\alpha$, $z = -c$ respectively. Show that their central axis lies on the line $y = \dfrac{x(P-Q)\tan\alpha}{P+Q}$ and $z = \dfrac{cP Q}{P^2 + 2PQ\cos 2\alpha + Q^2}$. \pts{6}
\end{parts}

\medskip\rule{\linewidth}{0.2pt}

\medskip\noindent\textbf{Question 3.}
\begin{parts}
  \item A solid homogeneous hemisphere of radius $r$ has a solid right cone of the same material constructed on its flat base; the hemisphere rests on the convex side of a fixed sphere of radius $R$, with the axis of the cone vertical. Show that the greatest height of the cone consistent with stability for a small rolling displacement is:
  \[
  h = \sqrt{\frac{(3R+r)(R-r)}{3}} - \frac{2r}{3}
  \] \pts{6}
  \item Two equal uniform rods are firmly jointed at one end so that the angle between them is $\alpha$, and they rest in a vertical plane on a smooth sphere of radius $r$. Show that they are in stable or unstable equilibrium according as the length of either rod is $\lessgtr 4r\,\text{cosec}\,\alpha$. \pts{6}
\end{parts}

\medskip\rule{\linewidth}{0.2pt}

\medskip\noindent\textbf{Question 4.}
\begin{parts}
  \item Obtain the equation of the momental ellipsoid of a rigid body at a point $O$. Hence show that at each point $O$ there exists a set of three mutually perpendicular axes such that the products of inertia about them (taken two at a time) all vanish. \pts{6}
  \item If $A$ and $B$ are the moments of inertia of a uniform lamina about two perpendicular axes $Ox$ and $Oy$ lying in its plane, and $F$ is the product of inertia of the lamina about these axes, show that the principal moments at $O$ are equal to $\dfrac{1}{2}\bigl(A+B \pm \sqrt{(A-B)^2 + 4F^2}\bigr)$. \pts{6}
\end{parts}

\medskip\rule{\linewidth}{0.2pt}

\medskip\noindent\textbf{Question 5.}
\begin{parts}
  \item State D'Alembert's principle and obtain the general equations of motion of a rigid body under the action of impulsive forces. \pts{6}
  \item A rod of length $2a$ is moving about one end with uniform angular velocity on a smooth horizontal plane. Suddenly this end is released and a point, distant $b$ from this end, is fixed. Find the motion, considering the cases when $b \lessgtr \dfrac{4a}{3}$. \pts{6}
\end{parts}

\medskip\rule{\linewidth}{0.2pt}

\medskip\noindent\textbf{Question 6.}
\begin{parts}
  \item A right cone of semi-vertical angle $\alpha$ can turn freely about an axis passing through the centre of its base and perpendicular to its axis. If the cone starts from rest with its axis horizontal, show that when the axis is vertical, the thrust on the fixed axis is to the weight of the cone as $(1 + 4\cos^2\alpha)$ to $(1 - \tfrac{1}{3}\cos^2\alpha)$. \pts{6}
  \item A weightless straight rod $ABC$ of length $2a$ is movable about the fixed end $A$ and carries two particles of equal mass, one at the midpoint $B$ and the other at end $C$. If the rod is held horizontal and then released, show that the angular velocity when vertical is $\sqrt{\dfrac{3g}{5a}}$ and that $\dfrac{5a}{3}$ is the length of the equivalent simple pendulum. \pts{6}
\end{parts}

\medskip\rule{\linewidth}{0.2pt}

\medskip\noindent\textbf{Question 7.}
A uniform sphere of radius $a$ rolls down an inclined plane rough enough to prevent any sliding. Discuss the motion and show that for pure rolling motion, $\mu > \dfrac{2}{7}\tan\alpha$, where $\alpha$ is the inclination of the plane to the horizontal and $\mu$ is the coefficient of friction. \pts{12}

\vfill\rule{\linewidth}{0.4pt}
\begin{center}\small
  \href{https://bhu-exam.vercel.app/}{Click here to visit our website}\\[4pt]
  \href{https://me-aryan.vercel.app/}{Click here to visit our contributor}\\[4pt]
  \href{https://quashan-me.vercel.app/}{Click here to visit our contributor}
\end{center}
\end{document}
